From Proposition~\ref{lem:ext} and~\eqref{eq:pb.4}, we deduce the following lower bound.
\begin{corollary}\label{pub:leading}
The following holds:
\begin{enumerate}[label=(\roman*)]
\item Assume that $\Gamma$ is $C^2$. There exists $C>0$ such that, for every $m\geq m_1$, $v\in H^1(\Omega)$,
\[Cm\|v\|_{H^1(\Omega)}^2\geq\Lambda_{m,m^{-1/2}}(v)\geq\left(m\|v\|_{L^2(\Gamma)}^2+\int_\Gamma\frac\kappa2|v|^2\dx\Gamma\right)-\frac Cm\|v\|_{L^2(\Gamma)}^2.\]
\item Assume that $\Gamma$ is $C^{2,1}$. There exists $C>0$ such that, for every $m\geq m_1$, $v\in H^1(\Omega)$,
\[\left(m\|v\|_{L^2(\Gamma)}^2+\int_\Gamma\frac\kappa2|v|^2\dx\Gamma\right)+o(1)\geq\Lambda_{m,m^{-1/2}}(v).\]
\end{enumerate}
Here, the term $o(1)$ depends on $v$ (not only on the $H^1$ norm of $v$).
\end{corollary}
\begin{proof}
By Proposition~\ref{lem:ext}, the lower bound of Point (i) follows. Let us focus on Point (ii).
By the extension theorem for Sobolev functions (see for instance \cite[\S5.4]{evans1998partial}), there exist a constant $C>0$ and, for all $v\in H^1(\Omega)$, a function $Ev\in H^1(\R^3)$ that extends $v$ and such that $\|Ev\|_{H^1(\R^3)}\leq C\|v\|_{H^1(\Omega)}$.
Let us define the test function $u_m$ by $u_m=v\widetilde u_m$ where
\[\widetilde u_m\circ\Phi(s,t)=\begin{cases}v_{m,\kappa(s),K(s)}(mt)&\text{for all }(s,t)\in\Gamma\times[0,m^{-1/2}],\\0&\text{for all }(s,t)\in\Gamma\times[m^{-1/2},+\infty).\end{cases}\]
Here, the function $v_m$ is defined in~\eqref{eq:formaltestfun1}.
Let us first prove a general formula. Consider $u\in H^2(\Omega';\R)$ and $v\in H^1(\Omega';\C^4)$. With an integration by parts and using the fact that $u$ is real-valued,
\[\begin{aligned}
\|\nabla(vu)\|_{L^2(\Omega')}^2&=\|u\nabla v+v\nabla u\|_{L^2(\Omega')}^2\\
&=\|u\nabla v\|_{L^2(\Omega')}^2+\|v\nabla u\|_{L^2(\Omega')}^2+2\Re\braket{u\nabla v,v\nabla u}_{\Omega'}\\
&=\|u\nabla v\|_{L^2(\Omega')}^2+\Re\braket{uv,-v\Delta u}_{\Omega'}-\Re\braket{v\partial_\n u,vu}_\Gamma\\
&=\|u\nabla v\|_{L^2(\Omega')}^2+\braket{uv,-v\Delta u}_{\Omega'}-\braket{v\partial_\n u,vu}_\Gamma.
\end{aligned}\]
With an integration by parts only in the tangential direction,
\[\braket{uv,-v\Delta u}_{\Omega'}=\braket{uv,-v\Delta_tu}_{\Omega'}+2\Re\braket{u\nabla_sv,v\nabla_su}_{\Omega'}+\|u\nabla_sv\|_{L^2(\Omega')}^2,\]
where $\nabla_s$ is the tangential derivative and $-\Delta_t$ is the part of the Laplacian involving the second order derivative in the normal variable $t$. Thus, we get
\[\begin{aligned}
\|\nabla(vu)\|_{L^2(\Omega')}^2={}&\|u\nabla v\|_{L^2(\Omega')}^2+\braket{uv,-v\Delta_tu}_{\Omega'}+2\Re\braket{u\nabla_sv,v\nabla_su}_{\Omega'}\\&+\|u\nabla_sv\|_{L^2(\Omega')}^2-\braket{v\partial_\n u,vu}_\Gamma.
\end{aligned}\]
By density, this formula can be extended to $u$ in $H_t^2$ and $H_s^1$. Therefore, we can replace $u$ by $\widetilde u_m$. We get
\begin{equation}\label{pub:2.19}\begin{aligned}
\mathcal{Q}_m(u_m)={}&-\braket{v\partial_\n\widetilde u_m,v\widetilde u_m}_\Gamma+\|\widetilde u_m\nabla v\|_{L^2(\Omega')}^2+\braket{\widetilde u_mv,v(-\Delta_t+m^2)\widetilde u_m}_{\Omega'}\\&+2\Re\braket{\widetilde u_m\nabla_sv,v\nabla_s\widetilde u_m}_{\Omega'}+\|\widetilde u_m\nabla_sv\|_{L^2(\Omega')}^2.
\end{aligned}\end{equation}
With the explicit expression~\eqref{eq:formaltestfun1}, we find
\begin{equation}\label{pub:2.20}-\braket{v\partial_\n\widetilde u_m,v\widetilde u_m}_\Gamma\leq m\|v\|_{L^2(\Gamma)}^2+\int_\Gamma\frac\kappa2|v|^2\dx\Gamma+\frac Cm\|v\|_{L^2(\Gamma)}^2.\end{equation}
By using the dominated convergence theorem and the explicit expression $\widetilde u_m$, we get that the other terms in~\eqref{pub:2.19} go to $0$. Note here that this argument uses at most one derivative of the functions $\kappa(\cdot)$ and $K(\cdot)$ (see the definition of $v_{m,k,K}$ in~\eqref{eq:formaltestfun1}). That is why we need $\Gamma$ to be $C^{2,1}$.
With the definition of $\Lambda_{m,m^{-1/2}}(v)$, we find
\[\Lambda_{m,m^{-1/2}}(v)\leq m\|v\|_{L^2(\Gamma)}^2+\int_\Gamma\frac\kappa2|v|^2\dx\Gamma+o(1).\]
To get the upper bound of Point (i), we follow the same steps as before except that $v_{m,k,K}$ is replaced by
\[\tau\longmapsto\chi_m(\tau)u_0(\tau),\]
in~\eqref{eq:formaltestfun1}. In that case, we only need $\Gamma$ to be $C^2$.
\end{proof}
Using Proposition~\ref{lem:agmon1}, Corollary~\ref{pub:leading} proves in particular (i) and (ii) in Proposition~\ref{prop:ext}. In this section we address the refinement of the lower bound and the corresponding upper bound. From now on, we assume that $\Gamma$ is $C^{2,1}$.
